\begin{titlepage}
\phantomsection
\label{sec:oc133-title-page}
\thispagestyle{empty}
\definecolor{ocTitleBlue}{HTML}{173B57}
\definecolor{ocTitleGold}{HTML}{A77D2A}
\begin{tikzpicture}[remember picture,overlay]
\draw[ocTitleBlue,line width=0.75pt]
  ([xshift=1.35cm,yshift=-1.35cm]current page.north west)
  rectangle
  ([xshift=-1.35cm,yshift=1.35cm]current page.south east);
\draw[ocTitleGold,line width=0.35pt]
  ([xshift=1.55cm,yshift=-1.55cm]current page.north west)
  rectangle
  ([xshift=-1.55cm,yshift=1.55cm]current page.south east);
\end{tikzpicture}
\begin{center}
\vspace*{1.0cm}
{\sffamily\Large\color{ocTitleBlue}\textsc{Ontology of Continua}}\\[0.38cm]
{\sffamily\Huge\bfseries Core~1.3.3}\\[0.28cm]
{\sffamily\Large Master Monograph}\\[0.75cm]
{\color{ocTitleGold}\rule{0.42\textwidth}{0.5pt}}\\[0.95cm]

{\Large Alexander Yashin}\\[0.2cm]
{\normalsize Independent Researcher, Leipzig/Halle, Germany}\\[0.2cm]
{\normalsize \href{https://orcid.org/0009-0008-6166-0914}{ORCID 0009-0008-6166-0914}}\\[1.0cm]

{\normalsize Version v1.3.3; revision recovery\_r005}\\[0.3cm]
{\normalsize Manuscript revision date: 3 May 2026}\\[0.3cm]
{\normalsize Concept DOI: \href{https://doi.org/10.5281/zenodo.17899134}{10.5281/zenodo.17899134}}\\[0.9cm]

\vfill

% DEDICATION_REQUIRED_DO_NOT_REMOVE
% OC_CORE_PUBLICATION_DEDICATION
{\small\itshape Dedicated to my dear wife Maria, without whom this work would have been impossible.}\\[0.8cm]

\begin{minipage}{0.86\textwidth}
\centering\scriptsize
This public monograph is the canonical OC Core~1.3.3 scientific manuscript.
Repository, archive, and evidence-package links are placed in the final reference section rather than on the title page.
\end{minipage}
\end{center}
\end{titlepage}

\clearpage
\noindent\textbf{Acknowledgements.}
The author thanks the reviewers and critics whose questions, objections, and
suggestions contributed to the development of the Ontology of Continua and to
the refinement of this release. Their criticism helped sharpen the claim
boundaries, improve the reader route, expose weak presentation choices, and
force clearer separation between model claims, proof routes, numerical evidence,
prior-art comparison, and publication form. Acknowledgement records review
pressure and intellectual contribution to the development of the work; it does
not imply authorship, endorsement, publication approval, responsibility for the
theory, or agreement with any claim promoted in the manuscript.
\begin{itemize}
    \item G.~V.~Apostolov.
    \item Eduard Fadeev.
    \item Gennady Alekseevich Nosov.
    \item Sergey Shpadyrev.
    \item Stanislav Tsukrov.
\end{itemize}

\clearpage
\begin{abstract}
Ontology of Continua (OC) Core~1.3.3 is the bounded external-review release of
the OC model core. It presents a typed account of continuants, realizations,
liveness, death, residue, rebirth, morphisms, generalized boundaries, hybrid
operators, cycle modes, historical and effective dimension, and K-level
witnesses under explicit assumptions. The scientific task of this manuscript is
not to replace every domain science with a slogan. Its task is to give reviewers
a precise object whose definitions, proof routes, finite semantic checks,
figures, tables, numerical anchors, comparator boundaries, and failure
conditions can be inspected.

The monograph is written as one continuous scientific manuscript rather than as
a prior volume followed by an update packet. The 1.3.3 additions are integrated
into the model, proof, evidence, comparison, limitation, and reviewer-boundary
chapters. The reader should therefore treat the text as the canonical public
spine of the release: it teaches the model first, then binds theorem routes,
Lean subset evidence, finite-model witnesses, bounded target-blind replay QA,
prior-art positioning, and adversarial-review material to the claim surface.

The current maturity of OC Core is deliberately stated as bounded. The release
is mature enough to be read as a coherent scientific manuscript and to be
checked against formal and executable artifacts. It is not a declaration of
complete scientific coverage, complete numerical closure, or unbounded
comparative victory over contemporary science. Where a theorem route, proof sheet,
Lean declaration, finite semantic case, replay row, comparator row, or
negative-control route supports a claim, the claim is promoted. Where that
support is incomplete, the text records a limitation or future obligation
instead of widening the public wording.

The manuscript contains the full monograph body, a compact article route, a
methods and reproducibility route, reviewer-response material, proof and
formalization sections, finite-model and numerical evidence anchors, figure and
table material, source-trace appendices, and practical-use boundaries. Figures
and tables are retained as scientific material; formulas and theorem statements
are retained in the mathematical spine; machine-readable registers remain in
the evidence package where they can be audited without becoming the narrative
structure of the manuscript.

The principal limitation is part of the method. A reader should evaluate whether
the typed model, proof governance, finite semantics, bounded replay evidence,
prior-art comparison, and adversarial-response route form a coherent model-core
contribution under stated assumptions. Future releases must add evidence,
mechanization, comparator work, or domain evidence-boundary review before they widen the
public claim surface. Version~1.3.3 therefore offers a publication-grade basis
for external review, not a claim that every future scientific obligation has
already been discharged.
\end{abstract}

\clearpage
\section*{Version 1.3.3 Release Delta}
\addcontentsline{toc}{section}{Version 1.3.3 Release Delta}
Version~1.3.3 is the release in which the OC Core corpus is reorganized from
scattered source witnesses and process records into a reviewable scientific
package. The visible delta is not merely a new archive number: the release
strengthens the typed foundation, makes the claim boundary explicit, integrates
the theorem route with public proof sheets, and separates reader-facing prose
from machine evidence.

The model delta includes clearer treatment of carriers, realizations, liveness,
death, residue, rebirth, morphisms, generalized boundaries, hybrid operators,
cycle modes, historical and effective dimension, and K-level witnesses. These
additions are presented as a bounded model-core grammar rather than as an
unrestricted theory of every phenomenon.

The verification delta adds and consolidates a Lean-checked subset, structured
proof sheets, finite semantic checks, finite witness interpretation, bounded
target-blind replay QA, comparator and prior-art positioning, negative-control
and falsifier language, and adversarial-review response material. Where a
stronger claim is not yet earned, the release records the limitation instead of
hiding it inside a source-control record.

The editorial delta is equally important. Figures, tables, formulas, numeric
anchors, appendices, and evidence routes are kept in the publication corpus;
machine coverage maps remain coverage and trace data only. The public
documents are therefore built from the human manuscript hierarchy and curated
payload sources, while source bindings and machine evidence indexes stay in the
review manifest.

\clearpage
\section*{Reader Routes}
\addcontentsline{toc}{section}{Reader Routes}
Dear reader, this monograph is the complete public manuscript for OC
Core~1.3.3. It is not an internal routing memo, not a build log, and not a
one-to-one printout of evidence-node records. It exists so that a careful
reader can understand the model, inspect the proof and evidence routes, and
then decide which machine-readable materials are worth opening.

\paragraph{Scientific reviewers and formal critics.}
Begin with the Abstract, Version Delta, the model-foundation chapters, and the
theorem/proof integration route. Then read the Lean subset, finite semantic
witnesses, proof machinery appendix, and falsifiability sections. Your main
question is whether each promoted theorem has assumptions, definitions,
dependency route, proof idea, executable witness where relevant, and a concrete
counterexample boundary.

\paragraph{Systems theorists, philosophers, and interested theoretical readers.}
Read the continuum problem, systems-theory context, K-level hierarchy, boundary
and lifecycle chapters, prior-art comparison, and limitations before opening
numeric evidence. Your route asks what OC borrows, what it reorganizes, where
its residual delta begins, and where existing systems theory, autopoiesis,
dynamical-systems, category-theoretic, or identity/persistence work already
carries part of the claim.

\paragraph{AI builders, engineers, architects, and applied researchers.}
Use the model chapters as a design grammar: carriers, realizations, boundaries,
operators, update/flow separation, state transitions, evidence classes, and
reopening rules. Then read the methods companion, worked examples, figures,
tables, reproducibility route, and practical-utility comparison sections. Your
route is useful when building AI systems, evaluation pipelines, knowledge
graphs, formalized domain models, or simulation protocols that need explicit
claim-to-evidence discipline.

\paragraph{CIO, CEO, enterprise architects, and strategy readers.}
Read the release delta, reader guide, practical-utility chapter, model
comparison appendix, and final conclusion first. Skip proof details on the
first pass unless a decision depends on them. Your route asks what the model is
for, what it can and cannot yet support, how it could influence enterprise
architecture, AI governance, evidence management, systems design, or research
portfolio strategy, and which claims remain too immature for operational
commitment.

\paragraph{Reproducibility auditors, evidence reviewers, and benchmark readers.}
Read the methods and evidence chapters, target-blind replay QA, numeric tables,
machine-readable table reference, audit trail, source-trace appendix, and public
evidence package manifest. Your route asks whether a claim has a formula or
reconstruction rule, input snapshot, comparator, residual or uncertainty,
negative control, falsifier, replay hash, and documented failure interpretation.

\paragraph{Reading rule shared by all groups.}
Read limitations as part of the claim, not as a footnote. A statement is
promoted only where its assumptions, proof or evidence class, comparator
boundary, formulas, appendices, numeric or formal anchor, figure/table
explanation, and reopening condition are visible. Claims about complete scientific coverage, complete
numerical closure, or unbounded comparative victory over contemporary science
are not promoted by this release.

\clearpage
